% =============================================================================
% Borrador tesis PCT — reemplazar preámbulo con plantilla UAN cuando esté lista
% =============================================================================
\documentclass[12pt,a4paper,openany]{report}

\usepackage[utf8]{inputenc}
\usepackage[T1]{fontenc}
\usepackage[spanish,es-tabla]{babel}
\usepackage{geometry}
\usepackage{graphicx}
\usepackage{booktabs}
\usepackage{longtable}
\usepackage{hyperref}
\usepackage{enumitem}
\usepackage{listings}
\usepackage{xcolor}
\usepackage{amsmath}
\usepackage{caption}
\usepackage{subcaption}

\geometry{left=3cm, right=2.5cm, top=3cm, bottom=2.5cm}

\hypersetup{
  colorlinks=true,
  linkcolor=blue!60!black,
  citecolor=green!50!black,
  urlcolor=blue!70!black,
}

\lstset{
  basicstyle=\ttfamily\small,
  breaklines=true,
  frame=single,
  backgroundcolor=\color{gray!8},
  keywordstyle=\color{blue!70!black},
  commentstyle=\color{green!40!black},
}

% --- Metadatos (ajustar con plantilla oficial) ---
\newcommand{\tesisTitulo}{Protocolo distribuido de control y reconfiguración topológica con continuidad lógica para redes multisalto sobre Wi-Fi convencional}
\newcommand{\tesisAutorA}{Juan Carlos Clavijo Triviño}
\newcommand{\tesisAutorB}{Brandon Stiven Ganzo Murcia}
\newcommand{\tesisUniversidad}{Universidad Antonio Nariño}
\newcommand{\tesisPrograma}{Ingeniería de Sistemas}
\newcommand{\tesisCiudad}{Bogotá D.C.}
\newcommand{\tesisAno}{2026}

\begin{document}

% --- Portada provisional ---
\begin{titlepage}
  \centering
  {\Large \tesisUniversidad\par}
  \vspace{0.5cm}
  {\large \tesisPrograma\par}
  \vspace{2cm}
  {\bfseries\LARGE \tesisTitulo\par}
  \vspace{1.5cm}
  {\large Protocolo de Control Topológico (PCT)\par}
  \vspace{2cm}
  {\large \tesisAutorA\\ \tesisAutorB\par}
  \vfill
  {\large \tesisCiudad\\ \tesisAno\par}
\end{titlepage}

\pagenumbering{roman}
\tableofcontents
\cleardoublepage

\pagenumbering{arabic}

% --- Capítulos (un archivo por capítulo) ---
\chapter{Introducción}
\label{cap:introduccion}

\section{Planteamiento del problema}
\label{sec:planteamiento}

Las redes de datos en entornos sin infraestructura fija ---desastres naturales, operaciones de campo, comunidades rurales o escenarios tácticos--- requieren conectividad entre dispositivos móviles sin depender de torres celulares ni puntos de acceso cableados. Las redes \emph{ad hoc} y \emph{mesh} permiten que cada nodo actúe como repetidor, formando topologías multisalto donde un mensaje puede atravesar varios dispositivos intermedios antes de llegar a su destino lógico.

En smartphones comerciales con Android, las opciones tradicionales presentan limitaciones severas: \emph{Wi-Fi Direct} fue concebido para enlaces punto a punto entre aplicaciones, no para construir árboles de decenas de nodos; modificar la capa MAC o inyectar tramas 802.11 exige \emph{root}; y el \emph{NAT} del kernel no ofrece enrutamiento lógico por identidad de aplicación. Un protocolo de capa de aplicación que emule funciones de red sobre Wi-Fi convencional, sin privilegios elevados, constituye una línea de investigación pertinente para ingeniería de sistemas.

Este proyecto de grado propone el \textbf{Protocolo de Control Topológico (PCT)}: un protocolo distribuido de control y reconfiguración topológica con continuidad lógica para redes multisalto sobre Wi-Fi convencional en dispositivos Android~12+ (API~31).

\section{Justificación}
\label{sec:justificacion}

La justificación del trabajo reposa en cuatro ejes:

\begin{enumerate}[label=\arabic*.]
  \item \textbf{Accesibilidad hardware.} Aprovechar smartphones comerciales con Wi-Fi Direct (mandatorio en CDD Android) evita hardware especializado.
  \item \textbf{Modelo homogéneo.} Exigir que todo nodo sea Group Owner (GO) de su propio grupo simplifica el arranque y la re-convergencia tras fallos parciales.
  \item \textbf{Identidad lógica estable.} Enrutar por UUID (\texttt{node\_id}) desacopla la topología lógica de las direcciones IPv4 locales, que cambian en cada subred SoftAP.
  \item \textbf{Continuidad de sesión.} Mantener \texttt{session\_epoch} ante reconexiones reduce la pérdida de contexto conversacional en aplicaciones de mensajería mesh.
\end{enumerate}

\section{Objetivo general}
\label{sec:objetivo-general}

Diseñar, implementar y validar un protocolo distribuido de control topológico (PCT) que permita formar y operar redes multisalto sobre Wi-Fi convencional en Android, con enrutamiento lógico por identidad, reconfiguración sin reinicio global y continuidad de sesión ante cambios de ruta.

\section{Objetivos específicos}
\label{sec:objetivos-especificos}

\begin{enumerate}[label=\arabic*.]
  \item Definir la arquitectura del protocolo PCT, incluyendo planos de descubrimiento, control y datos, con trazabilidad a requisitos funcionales y no funcionales.
  \item Especificar el modelo de enlace basado en GO homogéneo y asociación STA legacy al SoftAP del padre, prohibiendo \texttt{WifiP2pManager.connect()} para topología.
  \item Diseñar la capa de enlace (L2) con dos canales TCP por vecino (control :8765 y datos :8766) y registro formal de vecinos (\texttt{NeighborRegistry}).
  \item Diseñar la capa de red (L3) con tabla de rutas local, propagación \texttt{TOPO\_UPDATE} y reenvío multisalto de tráfico usuario.
  \item Implementar la biblioteca \texttt{lib-pct-core} (\texttt{co.uan.pct:core}) y una aplicación demo que evidencie bootstrap, topología y mensajería.
  \item Validar experimentalmente las hipótesis H1--H5 mediante la matriz EXP-01 a EXP-06 en laboratorio con dispositivos clase~$\alpha$.
\end{enumerate}

\section{Alcance y limitaciones}
\label{sec:alcance}

\subsection{Alcance}

\begin{itemize}
  \item Plataforma Android~12+ (\texttt{minSdk}~31), Kotlin, biblioteca Maven publicable.
  \item Topología física en árbol: un enlace STA upstream por nodo en el perfil mínimo.
  \item Profundidad máxima de siete saltos lógicos (\texttt{HOP\_LIMIT\_MAX}).
  \item Descubrimiento bootstrap vía DNS-SD P2P (\texttt{\_pct-ctrl}, \texttt{\_pct-seek}).
  \item Codificación wire binaria big-endian (v1).
\end{itemize}

\subsection{Limitaciones}

\begin{itemize}
  \item Sin cifrado extremo a extremo (RF-17, fuera de alcance v1).
  \item Sin QoS diferenciado (RF-18).
  \item PSK del GO en TXT DNS-SD solo para laboratorio controlado.
  \item Perfil hardware mínimo: clase~$\alpha$ (dual GO+STA concurrente verificado).
  \item Sin modificación de driver Wi-Fi ni acceso MAC/PHY.
\end{itemize}

\section{Estructura del documento}
\label{sec:estructura-documento}

El Capítulo~\ref{cap:marco-teorico} presenta fundamentos de redes mesh, Wi-Fi Direct y el modelo OSI emulado. El Capítulo~\ref{cap:marco-referencial} expone antecedentes, requisitos y restricciones Android. El Capítulo~\ref{cap:metodologia} describe la metodología incremental y la matriz experimental. El Capítulo~\ref{cap:analisis-diseno} detalla el análisis y diseño del sistema PCT. El Capítulo~\ref{cap:desarrollo} documenta la implementación en \texttt{lib-pct-core}. El Capítulo~\ref{cap:pruebas} reporta pruebas y resultados. El Capítulo~\ref{cap:conclusiones} cierra con conclusiones y trabajo futuro.

\chapter{Marco teórico y estado del arte}
\label{cap:marco-teorico}

\section{Redes \emph{ad hoc} y multisalto}
\label{sec:ad-hoc}

Una red \emph{ad hoc} móvil (MANET) es una colección de nodos inalámbricos que se comunican sin infraestructura preexistente~\cite{perkins2003adhoc}. Cada nodo puede actuar como host y como router. En topologías multisalto, un paquete atraviesa nodos intermedios antes de alcanzar el destino final.

Los protocolos de enrutamiento MANET se clasifican en proactivos (mantienen tablas actualizadas continuamente, p.\ ej.\ OLSR) y reactivos (descubren rutas bajo demanda, p.\ ej.\ AODV, DSDV)~\cite{rfc3561}. En entornos con recursos limitados de batería y CPU ---como smartphones--- el equilibrio entre overhead de control y latencia de convergencia es crítico.

\section{Redes mesh en dispositivos móviles}
\label{sec:mesh-moviles}

Soluciones comerciales y de investigación han explorado mesh sobre Android mediante Wi-Fi Direct, Bluetooth o tunelización. Proyectos como Serval Mesh y Commotion Wireless demostraron viabilidad en escenarios comunitarios, pero suelen requerir \emph{root}, firmware modificado o arquitecturas distintas al perfil homogéneo GO+STA que propone PCT.

PCT se distingue por operar \textbf{exclusivamente en espacio de usuario}, sin \texttt{connect()} P2P para topología, y por separar explícitamente identidad lógica (UUID) de dirección IPv4 local.

\section{Wi-Fi Direct y SoftAP}
\label{sec:wifi-direct}

Wi-Fi Direct (Wi-Fi P2P) permite que dispositivos negocien un grupo con un Group Owner (GO) que actúa como SoftAP~\cite{wifidirect2010}. Android expone la API \texttt{WifiP2pManager} con operaciones \texttt{createGroup()}, \texttt{removeGroup()} y \texttt{requestGroupInfo()}.

La asociación STA legacy (\texttt{WifiNetworkSpecifier} desde API~29) permite que un dispositivo se conecte al SSID/PSK de otro GO como estación infraestructura, habilitando el patrón ``doble interfaz lógica'': GO propio para hijos downstream y STA upstream hacia el padre.

\section{DNS-SD y descubrimiento de servicios}
\label{sec:dns-sd}

DNS-Based Service Discovery (DNS-SD) sobre mDNS permite anunciar y descubrir servicios en redes locales~\cite{rfc6763}. Android Wi-Fi P2P integra DNS-SD mediante \texttt{WifiP2pDnsSdServiceInfo} y \texttt{addLocalService()}. PCT utiliza registros TXT compactos para transportar \texttt{node\_id}, rol, SSID y PSK del GO en la fase bootstrap (capa L1).

\section{Modelo OSI como referencia conceptual}
\label{sec:osi-emulado}

PCT no implementa una pila OSI completa; utiliza el modelo como \textbf{mapa de responsabilidades}:

\begin{table}[htbp]
  \centering
  \caption{Correspondencia capas OSI -- PCT}
  \label{tab:osi-pct}
  \begin{tabular}{@{}lll@{}}
    \toprule
    \textbf{Capa OSI} & \textbf{Capa PCT} & \textbf{Transporte} \\
    \midrule
    Física (L1)       & Radio / bootstrap & 802.11, DNS-SD \\
    Enlace (L2)       & Vecindad          & TCP :8765, :8766 \\
    Red (L3)          & Reenvío lógico    & DATA en :8766 \\
    Aplicación        & Mensajería usuario & Envelope \texttt{type=user} \\
    \bottomrule
  \end{tabular}
\end{table}

La regla normativa es que la capa de red \textbf{nunca} enruta por IP de destino final; la IPv4 (\texttt{next\_hop\_local\_ip}) es un dato de enlace para alcanzar al vecino inmediato.

\section{Estado del arte y posicionamiento de PCT}
\label{sec:estado-arte}

% TODO: ampliar con revisión bibliográfica sistemática (Scopus/IEEE)
Protocolos como B.A.T.M.A.N.~\cite{openwrt_batman} operan en capa~2 con visión global de vecinos por broadcast periódico. Babel~\cite{rfc8966} es un protocolo de enrutamiento vector-distancia para IPv6. Ambos asumen control sobre la interfaz de red a bajo nivel.

PCT adopta un enfoque de \textbf{capa de aplicación}: construye topología y enrutamiento sobre TCP confiable vecino-a-vecino, con identidad UUID y continuidad lógica (\texttt{epoch}, \texttt{session\_epoch}). Esta decisión sacrifica integración con el stack IP del sistema operativo a cambio de despliegue sin privilegios en Android comercial.

\section{Síntesis del capítulo}
\label{sec:sintesis-marco-teorico}

El marco teórico establece que las MANET multisalto en móviles requieren equilibrar overhead, identidad estable y restricciones de plataforma. Wi-Fi Direct GO + STA legacy + DNS-SD constituyen la base física de PCT; el modelo OSI emulado organiza las responsabilidades L1--L3 que se diseñan e implementan en capítulos posteriores.

\chapter{Marco referencial y antecedentes}
\label{cap:marco-referencial}

\section{Antecedentes del proyecto}
\label{sec:antecedentes}

El presente trabajo de grado se desarrolla en la Universidad Antonio Nariño (UAN) como continuación del anteproyecto titulado \emph{Protocolo distribuido de control y reconfiguración topológica con continuidad lógica para redes multisalto sobre Wi-Fi convencional}. El anteproyecto estableció el planteamiento, objetivos preliminares y alcance; este documento profundiza en diseño formal, implementación y validación experimental.

Los autores ---Juan Carlos Clavijo Triviño y Brandon Stiven Ganzo Murcia--- dividieron responsabilidades según la arquitectura del sistema: conectividad Wi-Fi (GO, DNS-SD, STA, plano TCP) y enrutamiento/topología (tablas, heartbeats, reenvío).

\section{Prototipo EXP-01 (proto-exp01-gostat)}
\label{sec:exp01}

Como antecedente de implementación se desarrolló el prototipo experimental EXP-01, documentado en \texttt{docs/prototype/proto-exp01-gostat/}. Este prototipo demostró en laboratorio:

\begin{itemize}
  \item Formación de grupo P2P GO en dos dispositivos Android independientes.
  \item Descubrimiento DNS-SD del servicio \texttt{\_pct-ctrl} con parseo de TXT.
  \item Asociación STA legacy del hijo al SoftAP del padre.
  \item Coexistencia GO propio + STA upstream (hipótesis H1).
\end{itemize}

EXP-01 identificó restricciones operativas críticas: el buscador no debe mantener GO activo durante el escaneo DNS-SD; y la identidad de hijos no debe depender de \texttt{clientList} MAC cuando exista handshake formal.

\section{Evolución hacia lib-pct-core}
\label{sec:evolucion-core}

El prototipo EXP-01 se refactorizó en la biblioteca Maven \texttt{co.uan.pct:core} (\texttt{pct/library/lib-pct-core}), con API pública \texttt{PctNode}:

\begin{lstlisting}[language=Java,caption={API pública PCT (Kotlin)},label={lst:api-pct}]
val pct = PctCore.create()
pct.init(context, PctConfig.DEFAULT)
pct.start()  // bootstrap: scan -> join | root
pct.topology.collect { ... }
pct.sendUser(destinationNid, payload)
pct.close()
\end{lstlisting}

La biblioteca concentra bootstrap automático, teardown P2P anti-zombi al cerrar la aplicación, y observables (\texttt{StateFlow}) de fase, topología, vecinos y rutas.

\section{Requisitos funcionales}
\label{sec:requisitos-funcionales}

La Tabla~\ref{tab:rf} resume los requisitos funcionales trazables (RF). Los marcados \emph{Must} son obligatorios para la versión de grado.

\begin{longtable}{@{}clp{7cm}@{}}
  \caption{Requisitos funcionales PCT}
  \label{tab:rf} \\
  \toprule
  \textbf{ID} & \textbf{Pri.} & \textbf{Descripción} \\
  \midrule
  \endfirsthead
  \toprule
  \textbf{ID} & \textbf{Pri.} & \textbf{Descripción} \\
  \midrule
  \endhead
  RF-01 & Must & Identidad UUID única por nodo \\
  RF-02 & Must & Arranque homogéneo como P2P GO \\
  RF-03 & Must & Prohibición de \texttt{connect()} para topología \\
  RF-04 & Must & Enlace upstream STA legacy al padre \\
  RF-05 & Must & Anuncio DNS-SD en GO \\
  RF-06 & Must & Unión simétrica (\texttt{\_pct-seek} / \texttt{JOIN\_OFFER}) \\
  RF-07 & Must & Handshake TCP post-asociación (HELLO + JOIN\_COMMIT) \\
  RF-08 & Must & Tabla de rutas local por \texttt{node\_id} \\
  RF-09 & Must & Prevención de ciclos (\texttt{path\_trace}, \texttt{hop\_limit}) \\
  RF-10 & Must & Heartbeats PING/PONG \\
  RF-11 & Must & Detección de fallo y NODE\_DOWN \\
  RF-12 & Must & Reconfiguración sin reinicio global \\
  RF-13 & Should & Multisalto $\geq$ 3 nodos \\
  RF-14 & Should & Unión tardía sin reiniciar red existente \\
  RF-15 & Should & Continuidad \texttt{session\_epoch} \\
  RF-16 & Could & DAG lógico multipath \\
  RF-17 & Won't & Cifrado E2E (v1) \\
  RF-18 & Won't & QoS (v1) \\
  \bottomrule
\end{longtable}

\section{Requisitos no funcionales}
\label{sec:rnf}

\begin{itemize}
  \item[RNF-01] Plataforma mínima Android~12 (API~31).
  \item[RNF-02] Wi-Fi Direct mandatorio.
  \item[RNF-03] RAM mínima 2~GB.
  \item[RNF-04] Profundidad máxima 7 saltos.
  \item[RNF-05] Reconexión objetivo $\leq$ 15~s.
  \item[RNF-06] Puerto control TCP 8765 (configurable).
  \item[RNF-07] Codificación wire binaria big-endian.
  \item[RNF-08--09] Payload TCP $\leq$ 1400~B; mensaje usuario $\leq$ 1024~B.
\end{itemize}

\section{Restricciones de plataforma Android}
\label{sec:restricciones-android}

Android~12+ impone permisos \texttt{NEARBY\_WIFI\_DEVICES} (API~33+) y restricciones de descubrimiento P2P concurrente con GO activo. No es posible modificar tramas 802.11 ni implementar NAT multisalto en kernel sin \emph{root}. PCT opera íntegramente en espacio de usuario mediante \texttt{Network.bindSocket()} para seleccionar la interfaz STA upstream frente al GO propio.

\section{Perfiles de hardware}
\label{sec:perfiles-hardware}

Los dispositivos de laboratorio se clasifican en perfiles $\alpha$, $\beta$, $\gamma$ según capacidad dual GO+STA concurrente, RAM y batería. El perfil mínimo $\alpha$ exige verificación de hipótesis H1 (GO y STA simultáneos) antes de continuar experimentos multisalto.

\section{Síntesis del capítulo}
\label{sec:sintesis-marco-referencial}

El marco referencial vincula el anteproyecto, el prototipo EXP-01 y la biblioteca \texttt{lib-pct-core} con un conjunto trazable de requisitos y restricciones Android que condicionan todas las decisiones de diseño documentadas en el Capítulo~\ref{cap:analisis-diseno}.

\chapter{Metodología}
\label{cap:metodologia}

\section{Tipo de investigación}
\label{sec:tipo-investigacion}

El proyecto combina investigación aplicada y desarrollo tecnológico: a partir del planteamiento del anteproyecto se define un diseño formal del protocolo (Incremento~1), se implementa en software (Incremento~2 y siguientes) y se valida mediante experimentos controlados en laboratorio con dispositivos Android comerciales.

\section{Enfoque incremental}
\label{sec:enfoque-incremental}

La metodología sigue entregas incrementales trazables:

\begin{enumerate}[label=\textbf{I\arabic*.}]
  \item \textbf{Incremento 1 --- Diseño formal:} arquitectura, requisitos, wire format TCP, máquinas de estado, matriz de validación (\texttt{docs/protocol/}, versión spec 0.1.0-pre).
  \item \textbf{Incremento 2 --- Capas formales:} modelo OSI emulado, L2 (enlace TCP dual), L3 (red), envoltorio de paquetes (spec 0.2.1-draft).
  \item \textbf{Implementación por fases A--D:} enlace $\rightarrow$ tabla de rutas $\rightarrow$ reenvío multisalto $\rightarrow$ endurecimiento.
  \item \textbf{Validación experimental:} experimentos EXP-01 a EXP-06 con criterios de aceptación medibles.
\end{enumerate}

\section{Fases de implementación}
\label{sec:fases-implementacion}

La Tabla~\ref{tab:fases} resume las fases de la hoja de ruta (\texttt{17-implementation-roadmap.md}).

\begin{table}[htbp]
  \centering
  \caption{Fases de implementación lib-pct-core}
  \label{tab:fases}
  \begin{tabular}{@{}ll@{}}
    \toprule
    \textbf{Fase} & \textbf{Objetivo} \\
    \midrule
    A & Capa de enlace formal: TCP :8765/:8766, NeighborRegistry \\
    B & Tabla de rutas local y TOPO\_UPDATE \\
    C & Reenvío multisalto (ForwardWorker) en canal datos \\
    D & NODE\_DOWN, rutas STALE, persistencia UUID, \texttt{\_pct-seek} \\
    \bottomrule
  \end{tabular}
\end{table}

\section{Matriz de validación experimental}
\label{sec:matriz-validacion}

\subsection{Hipótesis}

\begin{table}[htbp]
  \centering
  \caption{Hipótesis principales}
  \label{tab:hipotesis}
  \begin{tabular}{@{}cl@{}}
    \toprule
    \textbf{ID} & \textbf{Hipótesis} \\
    \midrule
    H1 & GO propio + STA legacy upstream simultáneos en clase $\alpha$ \\
    H2 & Dos GO interconectados sin \texttt{connect()} \\
    H3 & Cadena de tres nodos con reenvío TCP \\
    H4 & Unión tardía sin reinicio de SSID \\
    H5 & Reconexión mantiene \texttt{session\_epoch} \\
    \bottomrule
  \end{tabular}
\end{table}

\subsection{Experimentos}

\textbf{EXP-01 (gate crítico):} interconexión de dos GO. Criterio: hijo mantiene GO propio tras STA; HELLO + JOIN\_COMMIT exitosos; $T_{\mathrm{join}} < 15$~s.

\textbf{EXP-02:} DNS-SD con GO local sin Group Client.

\textbf{EXP-03:} cadena ROOT $\leftarrow$ BRIDGE $\leftarrow$ LEAF; DATA de C a A vía reenvío en B.

\textbf{EXP-04:} unión tardía de nodo ISLAND a red operativa.

\textbf{EXP-05:} caída de BRIDGE y reconvergencia de sub-árbol.

\textbf{EXP-06:} LEAF anuncia \texttt{\_pct-seek} y recibe JOIN\_OFFER.

\section{Entorno de laboratorio}
\label{sec:entorno-lab}

\begin{itemize}
  \item \textbf{Dispositivos:} mínimo cuatro smartphones clase~$\alpha$ (p.\ ej.\ Samsung Galaxy A24, Vivo Y21s).
  \item \textbf{Herramientas:} Android Studio, Gradle, ADB, logcat (\texttt{PctMesh}).
  \item \textbf{Software:} app demo \texttt{pctmesh}, biblioteca \texttt{co.uan.pct:core}.
  \item \textbf{Red:} entorno indoor sin AP externo; grupos P2P aislados.
\end{itemize}

\section{Técnicas de recolección de datos}
\label{sec:recoleccion-datos}

\begin{enumerate}
  \item Logs estructurados en tag \texttt{PctMesh} (bootstrap, L2, DNS-SD).
  \item Snapshots de UI debug: fase, vecinos, rutas, diagnóstico DNS-SD.
  \item Medición de $T_{\mathrm{join}}$, $T_{\mathrm{reconfig}}$ con marca temporal en logs.
  \item Tests unitarios JUnit: codec TCP, parseo UUID, utilidades de red.
\end{enumerate}

\section{Síntesis del capítulo}
\label{sec:sintesis-metodologia}

La metodología garantiza trazabilidad desde requisitos hasta código y experimentos. EXP-01 actúa como compuerta: su fallo obliga a reclasificar hardware antes de pruebas multisalto.

\chapter{Análisis y diseño del sistema}
\label{cap:analisis-diseno}

\section{Arquitectura general}
\label{sec:arquitectura}

PCT se organiza en capas de presentación, protocolo y conectividad Wi-Fi (Figura~\ref{fig:arquitectura}). El \texttt{ProtocolOrchestrator} (\texttt{PctNodeImpl}) coordina JoinManager, LinkOrchestrator, RouteOrchestrator y TopologyManager.

% FIGURA: diagrama de componentes (exportar desde docs/protocol/03-architecture.md)
\begin{figure}[htbp]
  \centering
  \fbox{\parbox{0.85\textwidth}{\centering\vspace{2cm}
    \textit{[Figura pendiente: diagrama de componentes PCT]}\\
    \small Fuente: elaboración propia basada en \texttt{03-architecture.md}
  \vspace{2cm}}}
  \caption{Arquitectura de componentes PCT}
  \label{fig:arquitectura}
\end{figure}

\section{Modo de conexión: GO homogéneo + STA legacy}
\label{sec:modo-conexion}

\subsection{Principio rector}

Todo nodo ejecuta \texttt{createGroup()} y es GO de su propio grupo. Los enlaces del árbol se establecen \textbf{exclusivamente} mediante asociación STA legacy al SSID/PSK del GO padre. Queda prohibido \texttt{WifiP2pManager.connect()} para formar topología.

\subsection{Dos interfaces lógicas en un nodo BRIDGE}

\begin{table}[htbp]
  \centering
  \caption{Interfaces lógicas del nodo BRIDGE}
  \label{tab:interfaces-bridge}
  \begin{tabular}{@{}lll@{}}
    \toprule
    \textbf{Dirección} & \textbf{Objeto de red} & \textbf{Conoce} \\
    \midrule
    Upstream   & \texttt{Network} STA (\texttt{LegacyStaRepository}) & IP gateway del padre \\
    Downstream & GO P2P propio (\texttt{GoRepository})               & IPs DHCP de hijos \\
    \bottomrule
  \end{tabular}
\end{table}

El hijo obtiene la IP del padre como gateway de la red STA (\texttt{ConnectivityManager.getLinkProperties}). El padre obtiene la IP del hijo desde la dirección remota del socket TCP entrante en su GO (típicamente 192.168.49.x).

\subsection{Secuencia bootstrap miembro}

\begin{enumerate}
  \item Scan DNS-SD \texttt{\_pct-ctrl}; lectura TXT (\texttt{nid}, \texttt{go\_ssid}, \texttt{cp}).
  \item \texttt{requestNetwork} STA al padre.
  \item \texttt{createGroup()} --- GO BRIDGE propio.
  \item \texttt{addLocalService(\_pct-ctrl)}.
  \item TCP control al gateway STA :8765 --- HELLO, JOIN\_COMMIT.
  \item Negociación canal datos :8766 vía mensajes de control.
\end{enumerate}

\section{Capa de enlace (L2)}
\label{sec:diseno-l2}

\subsection{Dos canales TCP por vecino}

Android no garantiza \emph{fairness} en un único socket TCP mezclado. PCT exige:

\begin{itemize}
  \item \textbf{Control (:8765):} HELLO, JOIN\_COMMIT, PING/PONG, TOPO\_UPDATE, señalización DATA\_CHANNEL\_*.
  \item \textbf{Datos (:8766):} exclusivamente frames DATA (\texttt{type=user}).
\end{itemize}

\subsection{NeighborRegistry}

Cada entrada almacena: \texttt{neighbor\_nid}, \texttt{iface} (UPSTREAM/DOWNSTREAM), \texttt{local\_ip}, \texttt{ctrl\_socket}, \texttt{data\_socket}, \texttt{data\_channel\_state}, rol, hop, timestamps.

Invariantes: 0--1 upstream; 0--N downstream; identidad autoritativa = \texttt{sender\_nid} en HELLO.

\subsection{Secuencia L2}

\begin{enumerate}
  \item Hijo conecta TCP :8765 tras STA connected.
  \item HELLO bidireccional.
  \item JOIN\_COMMIT (hijo $\rightarrow$ padre).
  \item DATA\_CHANNEL\_OPEN (padre $\rightarrow$ hijo).
  \item Hijo conecta TCP :8766; DATA\_CHANNEL\_ACK.
  \item TOPO\_UPDATE inicial.
\end{enumerate}

\section{Capa de red (L3)}
\label{sec:diseno-l3}

\subsection{RoutingEntry}

Campos: \texttt{dest\_nid}, \texttt{next\_hop\_nid}, \texttt{hop\_count}, \texttt{path\_seq}, \texttt{status} (ACTIVE/BACKUP/STALE/BROKEN), \texttt{last\_seen\_ms}.

\subsection{Fusión TOPO\_UPDATE}

Al recibir entrada $(dest, hop, path\_seq, origin)$: aceptar si menor \texttt{hop\_count}; en empate, mayor \texttt{path\_seq}; rechazar si \texttt{origin == self}; decrementar TTL; marcar STALE tras timeout.

\subsection{ForwardWorker}

\begin{verbatim}
FUNCIÓN forward(msg DATA):
  SI msg.dst_nid == self: ENTREGAR a aplicación
  SI hop_limit == 0 O self in path_trace: DESCARTAR
  entry ← routing_table[msg.dst_nid]
  ENVIAR por data_socket(entry.next_hop_nid)  // solo :8766
\end{verbatim}

\section{Mensajes dirigidos y difusión}
\label{sec:mensajes-broadcast}

\subsection{Unicast}

HELLO, JOIN\_COMMIT, JOIN\_OFFER, DATA dirigido por UUID destino con reenvío hop-a-hop en canal :8766.

\subsection{Difusión acotada}

\texttt{TOPO\_UPDATE} y \texttt{NODE\_DOWN} se propagan a vecinos directos con TTL (\texttt{TOPO\_TTL\_DEFAULT = 7}), sin flooding global sin control. DNS-SD actúa como broadcast de bootstrap en L1, no como canal de datos multisalto.

\section{Desconexiones en BRIDGE}
\label{sec:diseno-desconexion-bridge}

Detección: timeout PING/PONG, \texttt{Network.onLost} STA, EOF en socket upstream.

Secuencia:
\begin{enumerate}
  \item Cerrar sockets upstream; purgar NeighborRegistry y rutas vía padre.
  \item Emitir NODE\_DOWN a hijos downstream.
  \item Transitar a ISLAND; mantener GO; anunciar \texttt{\_pct-seek}.
  \item Re-JOIN; \texttt{session\_epoch} sin cambio; \texttt{epoch} red solo si ROOT caído.
\end{enumerate}

Caída solo del canal :8766: DATA\_CHANNEL\_RESET vía control, sin reiniciar GO/STA.

\section{Unión de árboles independientes}
\label{sec:union-arboles}

Dos sub-redes con ROOT y \texttt{epoch} distintos pueden fusionarse cuando un ISLAND de árbol~B recibe JOIN\_OFFER de un BRIDGE de árbol~A. El árbol absorbente incrementa \texttt{tree\_version}; las rutas cruzadas se aprenden por TOPO\_UPDATE. Anti-ciclos: \texttt{path\_trace} + \texttt{hop\_limit}.

\section{Servicios DNS-SD}
\label{sec:dns-sd-diseno}

\begin{table}[htbp]
  \centering
  \caption{Servicios DNS-SD PCT}
  \label{tab:dns-sd}
  \begin{tabular}{@{}lll@{}}
    \toprule
    \textbf{Servicio} & \textbf{Anunciante} & \textbf{Propósito} \\
    \midrule
    \texttt{\_pct-seek.\_tcp}  & ISLAND  & Intención de unión \\
    \texttt{\_pct-ctrl.\_tcp}  & ROOT, BRIDGE, LEAF & Control operativo \\
    \bottomrule
  \end{tabular}
\end{table}

\section{Máquinas de estado}
\label{sec:maquinas-estado}

Estados radio: \texttt{R\_BOOT}, \texttt{R\_GO\_READY}, \texttt{R\_GO\_UPSTREAM}, \texttt{R\_ISLAND\_SEEK}. Roles lógicos: ISLAND, ROOT, BRIDGE, LEAF. Transición crítica: pérdida upstream $\rightarrow$ ISLAND\_SEEK $\rightarrow$ re-JOIN $\rightarrow$ GO\_UPSTREAM.

\section{Síntesis del capítulo}
\label{sec:sintesis-analisis-diseno}

El diseño de PCT separa bootstrap (L1), vecindad confiable (L2) y reenvío lógico (L3), con dos objetos de red complementarios: STA para el padre y GO P2P para hijos. Las decisiones normativas prohiben atajos (connect P2P, UDP crítico, mezcla control/datos) que comprometan robustez en Android.

\chapter{Desarrollo e implementación}
\label{cap:desarrollo}

\section{Stack tecnológico}
\label{sec:stack}

\begin{itemize}
  \item Lenguaje: Kotlin 1.9+
  \item Build: Gradle (Kotlin DSL), Android Library Module
  \item \texttt{minSdk} 31, \texttt{targetSdk} 36
  \item Coroutines + \texttt{StateFlow} para concurrencia reactiva
  \item Publicación: Maven Local (\texttt{co.uan.pct:core})
  \item App demo: Jetpack Compose (\texttt{pct/app/app-demo})
\end{itemize}

\section{Estructura de lib-pct-core}
\label{sec:estructura-core}

\begin{table}[htbp]
  \centering
  \caption{Paquetes principales de lib-pct-core}
  \label{tab:paquetes}
  \begin{tabular}{@{}ll@{}}
    \toprule
    \textbf{Paquete} & \textbf{Responsabilidad} \\
    \midrule
    \texttt{api/} & \texttt{PctNode}, \texttt{PctConfig}, snapshots públicos \\
    \texttt{internal/p2p/} & \texttt{GoRepository}, \texttt{DnsSdRepository}, \texttt{P2pChannelHolder} \\
    \texttt{internal/sta/} & \texttt{LegacyStaRepository}, \texttt{WifiNetworkSpecifier} \\
    \texttt{internal/link/} & \texttt{LinkOrchestrator}, \texttt{NeighborRegistry} \\
    \texttt{internal/route/} & \texttt{RouteOrchestrator}, tabla y reenvío \\
    \texttt{internal/tcp/} & \texttt{PctFrameCodec}, sockets :8765/:8766 \\
    \texttt{internal/util/} & \texttt{PctNid}, \texttt{ParentSelector}, \texttt{PctNetworkHelper} \\
    \bottomrule
  \end{tabular}
\end{table}

\section{Capa física y bootstrap (L1)}
\label{sec:impl-l1}

\subsection{P2pChannelHolder}

Centraliza \texttt{WifiP2pManager.Channel}, \texttt{BroadcastReceiver} para eventos P2P y \texttt{forceTeardownP2p()} al cerrar la app (anti-proceso zombi).

\subsection{GoRepository}

Implementa \texttt{createGroup()}, \texttt{requestGroupInfo()} con reintentos, mapeo de \texttt{clientList} a \texttt{GoClient}, estados \texttt{GoState.Ready|Creating|Error}.

\subsection{LegacyStaRepository}

Mantiene \texttt{activeNetwork: Network?} --- la red STA al SoftAP del padre. Usa \texttt{ConnectivityManager.requestNetwork} con \texttt{WifiNetworkSpecifier}; el socket upstream usa \texttt{network.bindSocket()} explícito (no \texttt{bindProcessToNetwork} global).

\subsection{DnsSdRepository}

Registra \texttt{\_pct-ctrl} con TXT compacto (\texttt{n}, rol, hop, SSID, PSK). Descubrimiento con \texttt{discoverServices}; merge de candidatos instance+TXT; selección de padre vía \texttt{ParentSelector}.

\subsection{Bootstrap automático (PctNodeImpl)}

Secuencia en \texttt{bootstrap()}:
\begin{enumerate}
  \item \texttt{removeGroup()} si hay GO activo (requisito Android para escaneo fiable).
  \item Scan DNS-SD con \texttt{scanSettleMs}.
  \item Si hay candidato: STA $\rightarrow$ \texttt{createGroup()} BRIDGE $\rightarrow$ anuncio $\rightarrow$ L2 upstream.
  \item Si no: \texttt{startAsRoot()}.
\end{enumerate}

\section{Capa de enlace (L2)}
\label{sec:impl-l2}

\subsection{PctFrameCodec}

Framing \texttt{PCT1}: cabecera 12~bytes (\texttt{msg\_type}, longitud payload), payloads binarios para HELLO (59~B), JOIN\_COMMIT (40~B), PING/PONG (32~B), DATA, TOPO\_UPDATE.

\subsection{PctLinkSocket}

Clase base con \texttt{connectOutbound(host, port, network)}: crea \texttt{Socket}, \texttt{network.bindSocket()}, \texttt{connect()}. Evita self-connect en BRIDGE cuando GO propio también es 192.168.49.1.

\subsection{LinkOrchestrator}

\begin{itemize}
  \item \texttt{onGoReady()}: lanza \texttt{ControlAcceptWorker} y \texttt{DataAcceptWorker}.
  \item \texttt{onStaConnected()}: conecta upstream a gateway STA; HELLO; JOIN\_COMMIT.
  \item \texttt{NeighborRegistry}: upsert/remove/update thread-safe.
  \item \texttt{canonicalParentNid} desde HELLO wire (UUID completo 128-bit).
  \item Recuperación upstream en EOF; DATA\_CHANNEL\_OPEN/ACK/RESET.
\end{itemize}

\subsection{Identidad UUID (PctNid)}

Rechaza UUID falsos con padding de ceros; merge instance DNS (8 hex) + TXT (\texttt{n=32 hex}); validación \texttt{isFull()}.

\section{Capa de red (L3)}
\label{sec:impl-l3}

\subsection{RouteOrchestrator}

Construye rutas directas desde \texttt{NeighborRegistry}; procesa TOPO\_UPDATE; expone \texttt{RouteSnapshot} y \texttt{NeighborSnapshot} vía \texttt{StateFlow}.

\subsection{ForwardWorker y sendUser}

\texttt{sendUser(destinationNid, payload)} encapsula \texttt{UserDataPayload} y delega reenvío. Entrega local si \texttt{dst == self}; decrementa \texttt{hop\_limit}; verifica \texttt{path\_trace}.

\section{Aplicación demo PCT Mesh}
\label{sec:app-demo}

La app \texttt{co.uan.pct.app.pctmesh} integra la biblioteca con:
\begin{itemize}
  \item Pantalla Debug: fase, topología, vecinos, rutas, log \texttt{PctMesh}.
  \item Pantalla Chat: \texttt{sendUser} entre nodos de la mesh.
  \item Permisos: \texttt{NEARBY\_WIFI\_DEVICES}, \texttt{ACCESS\_FINE\_LOCATION} ($\leq$ API~32).
\end{itemize}

\section{Convenciones de empaquetado UAN}
\label{sec:empaquetado}

Namespace \texttt{co.uan.pct}; experimentos numerados EXP-01..06; documentación en \texttt{docs/protocol/}; publicación local Gradle \texttt{publishToMavenLocal}.

\section{Síntesis del capítulo}
\label{sec:sintesis-desarrollo}

La implementación materializa el diseño en módulos Kotlin desacoplados. La separación \texttt{LegacyStaRepository.activeNetwork} (padre) vs aceptación TCP en GO (hijos) resuelve la ambigüedad de direccionamiento en nodos BRIDGE dual-interfaz.

\chapter{Pruebas, validación y resultados}
\label{cap:pruebas}

\section{Plan de pruebas}
\label{sec:plan-pruebas}

Las pruebas se organizan en tres niveles:

\begin{enumerate}
  \item \textbf{Unitarias:} codecs TCP, parseo UUID/NID, utilidades de gateway IP (\texttt{PctFrameCodecTest}, \texttt{PctNidTest}, \texttt{PctNetworkHelperTest}).
  \item \textbf{Integración L2:} handshake HELLO/JOIN en dos dispositivos; verificación UUID en topología.
  \item \textbf{Sistema (EXP-01..06):} escenarios multisalto en laboratorio con ADB/logcat.
\end{enumerate}

\section{EXP-01: Interconexión dos GO}
\label{sec:resultados-exp01}

\textbf{Configuración:} D1 ROOT (Samsung A24), D2 hijo (Vivo Y21s).

\textbf{Procedimiento:} D1 \texttt{start()} $\rightarrow$ espera anuncio \texttt{\_pct-ctrl} $\rightarrow$ D2 \texttt{start()}.

\textbf{Criterios:}
\begin{itemize}
  \item D2 asocia STA al SSID de D1.
  \item D2 mantiene GO propio (H1).
  \item TCP HELLO upstream completado; UUID padre visible en UI.
  \item $T_{\mathrm{join}} < 15$~s.
\end{itemize}

% TODO: completar tabla de resultados medidos en defensa
\begin{table}[htbp]
  \centering
  \caption{Resultados EXP-01 (plantilla --- completar con mediciones)}
  \label{tab:exp01}
  \begin{tabular}{@{}lcc@{}}
    \toprule
    \textbf{Métrica} & \textbf{Objetivo} & \textbf{Observado} \\
    \midrule
    STA asociada            & Sí  & --- \\
    GO hijo activo post-STA & Sí  & --- \\
    HELLO L2 completado     & Sí  & --- \\
    $T_{\mathrm{join}}$ (ms) & $<15000$ & --- \\
    \bottomrule
  \end{tabular}
\end{table}

\section{EXP-03: Cadena tres nodos}
\label{sec:resultados-exp03}

Topología: A (ROOT) $\leftarrow$ B (BRIDGE) $\leftarrow$ C (LEAF). C envía DATA a A; B reenvía por \texttt{ForwardWorker}. Criterio: entrega en C con \texttt{hop\_limit} decrementado y ruta ACTIVE en tablas intermedias.

\section{EXP-04 y EXP-05: Unión tardía y fallo BRIDGE}
\label{sec:exp04-exp05}

EXP-04 verifica que un nodo se una a red operativa sin \texttt{removeGroup} en nodos existentes. EXP-05 simula caída de B: hijos de B detectan NODE\_DOWN; B re-converge vía \texttt{\_pct-seek}; \texttt{session\_epoch} de conversaciones activas se mantiene (H5).

\section{Pruebas unitarias}
\label{sec:pruebas-unitarias}

Ejecución:
\begin{verbatim}
cd pct/library && ./gradlew :lib-pct-core:test
\end{verbatim}

Las pruebas validan round-trip de frames PCT1, normalización de UUID y resolución de gateway IPv4 en \texttt{Network} STA.

\section{Análisis de logs}
\label{sec:analisis-logs}

Tag \texttt{PctMesh} registra: fases bootstrap, candidatos DNS-SD, eventos STA, conexiones TCP (\texttt{local=}/\texttt{remote=}), HELLO upstream/downstream, errores self-connect.

Indicadores de fallo diagnosticados:
\begin{itemize}
  \item \texttt{pct=0 candidatos=0} con peers visibles: hijo escaneó antes del anuncio ROOT.
  \item \texttt{local=192.168.49.1 remote=192.168.49.1}: bind a GO propio en lugar de STA.
  \item UUID truncado con ceros: instance DNS sin merge TXT.
\end{itemize}

\section{Discusión}
\label{sec:discusion-resultados}

Los resultados confirman viabilidad de H1/H2 en hardware clase~$\alpha$. La validación multisalto (H3) depende de L3 estable y rutas coherentes tras TOPO\_UPDATE. Las métricas de reconexión (H5) requieren escenario EXP-05 repetido con carga de mensajería activa.

\section{Síntesis del capítulo}
\label{sec:sintesis-pruebas}

La matriz experimental proporciona criterios objetivos de aceptación. Las plantillas de resultados se completan con mediciones de la sesión de defensa; la infraestructura de prueba (ADB, logs, tests unitarios) ya está integrada en el repositorio.

\chapter{Conclusiones y trabajo futuro}
\label{cap:conclusiones}

\section{Conclusiones}
\label{sec:conclusiones}

\begin{enumerate}
  \item Se diseñó e implementó el protocolo PCT, un sistema distribuido de control topológico para redes multisalto sobre Wi-Fi convencional en Android~12+, cumpliendo la restricción de operar sin \emph{root} ni modificación del driver inalámbrico.

  \item El modelo de enlace GO homogéneo + STA legacy demostró ser viable en dispositivos clase~$\alpha$, validando las hipótesis H1 y H2 cuando el bootstrap respeta la separación entre red STA (padre) y GO P2P (hijos).

  \item La arquitectura por capas emuladas (L1 bootstrap, L2 vecindad TCP dual, L3 reenvío UUID) permite evolucionar el sistema de forma incremental sin acoplar descubrimiento DNS-SD al plano de datos multisalto.

  \item La identidad lógica basada en UUID de 128 bits, intercambiada en HELLO TCP, resuelve la ambigüedad de direcciones MAC en \texttt{clientList} y habilita trazabilidad en UI y enrutamiento.

  \item La separación de canales TCP :8765 (control) y :8766 (datos) es una decisión de diseño necesaria en Android para evitar que tráfico usuario bloquee heartbeats y sincronización topológica.

  \item La biblioteca \texttt{co.uan.pct:core} y la aplicación demo PCT Mesh constituyen un artefacto reproducible, publicable vía Maven Local e integrable en proyectos de mensajería off-grid.
\end{enumerate}

\section{Cumplimiento de objetivos}
\label{sec:cumplimiento-objetivos}

\begin{table}[htbp]
  \centering
  \caption{Trazabilidad objetivos específicos}
  \label{tab:objetivos-cumplimiento}
  \begin{tabular}{@{}clc@{}}
    \toprule
    \textbf{Obj.} & \textbf{Descripción} & \textbf{Estado} \\
    \midrule
    1 & Arquitectura y requisitos & Completado (spec + tesis) \\
    2 & Modelo GO + STA legacy     & Completado \\
    3 & Capa L2 dual TCP           & Implementado \\
    4 & Capa L3 reenvío            & Implementado / en validación \\
    5 & lib-pct-core + app demo    & Completado \\
    6 & Validación EXP-01..06      & En curso \\
    \bottomrule
  \end{tabular}
\end{table}

\section{Limitaciones del trabajo}
\label{sec:limitaciones-trabajo}

\begin{itemize}
  \item PSK en TXT DNS-SD expuesto en laboratorio; no apto para producción.
  \item Profundidad y tamaño de red limitados (7 hops, 32 rutas).
  \item Dependencia de fabricante en coexistencia GO+STA.
  \item Sin cifrado extremo a extremo ni QoS (RF-17, RF-18).
  \item Validación en entorno indoor controlado; no se evaluó movilidad a velocidad vehicular.
\end{itemize}

\section{Trabajo futuro}
\label{sec:trabajo-futuro}

\begin{enumerate}
  \item \textbf{Seguridad:} intercambio de claves fuera de TXT; cifrado E2E sobre canal :8766.
  \item \textbf{Persistencia:} almacenar \texttt{node\_id} en almacenamiento seguro; sobrevivir reinstalaciones.
  \item \textbf{DAG multipath (RF-16):} rutas BACKUP y selección de next-hop.
  \item \textbf{Dual-STA:} perfil extendido con \texttt{CAP\_DUAL\_STA}.
  \item \textbf{Publicación Maven} en repositorio institucional o GitHub Packages.
  \item \textbf{Pruebas de campo:} escenarios outdoor, movilidad, interferencia.
  \item \textbf{Integración IoT:} gateway PCT hacia infraestructura IP fija.
\end{enumerate}

\section{Cierre}
\label{sec:cierre}

PCT demuestra que es posible construir redes mesh lógicas sobre smartphones Android comerciales, usando Wi-Fi Direct y STA legacy como primitivas físicas y un protocolo de aplicación con identidad UUID y continuidad de sesión. El trabajo sienta bases formales (especificación en \texttt{docs/protocol/}), implementación (\texttt{lib-pct-core}) y metodología de validación reproducible para extensiones futuras.


\bibliographystyle{plain}
\bibliography{referencias/bibliografia}

\end{document}
